\documentclass[12pt, a4paper]{article}

%%------------------------------------------------------------------
%% Packages
%%------------------------------------------------------------------
\usepackage[T1]{fontenc}
\usepackage[utf8]{inputenc}
\usepackage{lmodern}
\usepackage{microtype}
\usepackage{amsmath, amssymb}
\usepackage{geometry}
\geometry{top=2.5cm, bottom=2.5cm, left=2.5cm, right=2.5cm}
\usepackage{graphicx}
\usepackage{booktabs}
\usepackage{array}
\usepackage{longtable}
\usepackage{caption}
\usepackage{hyperref}
\hypersetup{
  colorlinks = true,
  linkcolor  = blue,
  citecolor  = blue,
  urlcolor   = blue
}
\usepackage{natbib}
\bibliographystyle{unsrtnat}
\usepackage{parskip}
\setlength{\parindent}{0pt}
\usepackage{enumitem}
\usepackage{xcolor}
\usepackage{titlesec}
\titleformat{\section}{\large\bfseries}{\thesection}{1em}{}
\titleformat{\subsection}{\normalsize\bfseries}{\thesubsection}{1em}{}

%%------------------------------------------------------------------
%% Metadata
%%------------------------------------------------------------------
\title{\textbf{Assessing the Quality of Collective Variables in\\
Molecular Dynamics Enhanced Sampling:\\
A Review of Current Metrics and Methods}}
\author{Prepared for: B.~Bharland\\[6pt]
\small Computational Biophysics / Enhanced Sampling Group}
\date{June 2026}

\begin{document}
\maketitle
\tableofcontents
\newpage

%%==================================================================
\section{Introduction}
%%==================================================================

Enhanced sampling methods such as metadynamics~\citep{Laio2002}, 
well-tempered metadynamics~\citep{Barducci2008}, 
umbrella sampling, and replica exchange rely critically on a small set 
of \emph{collective variables} (CVs, also called order parameters or 
reaction coordinates) that encode the slow, large-scale motions 
relevant to the physical or chemical process under study.  When CVs are 
learned from simulation data---for instance by training a neural network 
on trajectory data---an immediate and practically important question 
arises: \textit{how good are the resulting CVs?}

This question has no single universal answer, because ``quality'' can be 
defined from several complementary perspectives:
\begin{enumerate}[noitemsep]
  \item \textbf{Kinetic quality}: Do the CVs correctly capture the slowest 
        dynamical modes, maximising timescale separation?
  \item \textbf{Thermodynamic quality}: Does biasing along the CVs lead to 
        rapid, converged estimates of the free energy?
  \item \textbf{Mechanistic quality}: Do the CVs identify the transition 
        state ensemble and the dominant reaction mechanism?
  \item \textbf{Informational quality}: Do the CVs retain maximum mutual 
        information about the microscopic slow coordinates?
\end{enumerate}

This report reviews the state of the art (roughly 2016--2025) in CV 
quality assessment, organised by the dominant methodological paradigm.  
Section~\ref{sec:committor} introduces the committor function as the 
theoretical gold standard.  Sections~\ref{sec:kinetic} 
and~\ref{sec:correlation} cover kinetic-based and time-correlation-based 
metrics, respectively.  Section~\ref{sec:vamp} discusses variational 
approaches and the VAMP score.  Sections~\ref{sec:freeenergy} 
and~\ref{sec:infotheory} cover free-energy convergence and 
information-theoretic metrics.  Section~\ref{sec:ml} summarises 
machine-learning frameworks that incorporate quality metrics by 
construction.  Section~\ref{sec:summary} provides an overall comparison 
and practical recommendations.

%%==================================================================
\section{The Theoretical Gold Standard: The Committor Function}
\label{sec:committor}
%%==================================================================

\subsection{Definition and properties}

For a system characterised by two metastable states $A$ (reactant) and 
$B$ (product), the \emph{committor function} $q(\mathbf{x})$ is defined 
as the probability that a trajectory initiated at configuration 
$\mathbf{x}$ with momenta drawn from the Boltzmann distribution will 
reach $B$ before $A$~\citep{EWainanBerkovitz2004}:
\begin{equation}
  q(\mathbf{x}) = P\bigl[\text{reach }B \text{ before }A \mid \mathbf{x}\bigr].
\end{equation}
The committor satisfies $q(\mathbf{x})=0$ for $\mathbf{x}\in A$ and 
$q(\mathbf{x})=1$ for $\mathbf{x}\in B$, and it is the solution to 
the backward Kolmogorov (Fokker--Planck) equation in the transition 
region.  Crucially, the committor is the \emph{optimal one-dimensional 
reaction coordinate} in the sense that its level sets ($q=\text{const}$) 
are the isocommittor (isoreactive) surfaces---the natural dividing 
surfaces of transition state theory.

Any candidate CV $s(\mathbf{x})$ can therefore be assessed by asking how 
well its level sets approximate the isocommittor surfaces 
of the true committor.

\subsection{The committor histogram test}

The classical operational test~\citep{Du1998,Geissler1999} is:
\begin{enumerate}[noitemsep]
  \item Identify configurations at the putative dividing surface 
        (where $s=s^*$, nominally the transition state).
  \item Launch short unbiased MD trajectories from each such 
        configuration with random momenta.
  \item Record whether each trajectory reaches $A$ or $B$ first and 
        compute the empirical committor $\hat{q}$.
  \item \textbf{Quality criterion}: If $\hat{q}$ is narrowly distributed 
        around $0.5$, the CV is good; a broad or bimodal distribution 
        signals that important slow coordinates are missing.
\end{enumerate}
This test is computationally expensive but remains the most direct 
validation available.

\subsection{Machine-learning approaches to computing the committor}

Recent work has made committor-based validation far more practical by 
using neural networks to \emph{learn} the committor directly.  Notably:

\begin{itemize}[noitemsep]
  \item \textbf{Committor-consistent variational string method}~\citep{Li2024CCVSM}: 
        Uses iterative committor computation via variational principles to 
        refine the minimum free-energy path and the transition state 
        ensemble simultaneously.
  \item \textbf{Committor-regularised learning of differentiable CVs} 
        (ChemRxiv 2025,~\citealt{Bonati2025CommittorReg}): Introduces a 
        regularisation term that penalises deviations of the learned CV 
        from committor-like behaviour, directly combining 
        interpretability (via physically meaningful descriptors) with 
        committor accuracy.
  \item \textbf{Self-consistent iterative committor sampling}~\citep{Kang2024}: 
        Bootstraps the committor from a current estimate and uses it to 
        guide trajectory sampling, achieving efficient TSE sampling and 
        rate computation.
  \item \textbf{Committor flow (data-driven, 2025)}~\citep{CommittorFlow2025}: 
        Uses normalising flows to parametrise the committor, enabling 
        data-efficient learning of transition mechanism and quality 
        assessment from short trajectory fragments.
\end{itemize}

The \texttt{mlcolvar} library~\citep{Bonati2023mlcolvar} provides a 
practical tutorial and implementation for learning the committor as a 
CV directly from labelled trajectory data.

%%==================================================================
\section{Kinetic-Based Quality Metrics}
\label{sec:kinetic}
%%==================================================================

\subsection{Barrier crossing rates as a direct quality indicator}

Perhaps the most physically intuitive test of a CV is whether it 
correctly captures barrier-crossing kinetics.  If the projected dynamics 
along the CV reproduce the true transition rates---as estimated from long 
unbiased simulations, transition interface sampling (TIS), or forward 
flux sampling (FFS)---the CV is considered high quality.  

In the context of metadynamics, reweighted kinetics can be extracted 
via the \emph{infrequent metadynamics} protocol~\citep{Tiwary2015IM}:
\begin{enumerate}[noitemsep]
  \item Deposit bias slowly (large deposition time $\tau_G$) so that the 
        barrier is rarely filled and individual transition times $t_i$ 
        can be identified.
  \item Apply the reweighting factor $e^{V(s,t)/k_BT}$ to correct each 
        observed transition time.
  \item Fit the corrected times to a Poisson distribution.
  \item \textbf{Quality criterion}: Passing a Kolmogorov--Smirnov (KS) 
        test against the ideal exponential distribution confirms that the 
        bias is deposited quasi-statically, i.e.\ that the CV captures 
        all relevant slow modes and no unphysical acceleration has 
        occurred.
\end{enumerate}
A failed KS test is an unambiguous signal that the CV set is 
\emph{incomplete}---important slow degrees of freedom are not biased 
and ``hidden'' recrossings corrupt the kinetics.

\subsection{Spectral Gap Optimization of Order Parameters (SGOOP)}
\label{sec:sgoop}

SGOOP~\citep{Tiwary2016SGOOP} formulates CV quality as a 
\emph{spectral gap maximisation} problem.  The key objects are:
\begin{enumerate}[noitemsep]
  \item A \emph{dictionary} of candidate order parameters 
        $\{\phi_i(\mathbf{x})\}$.
  \item A one-dimensional projected dynamics on a linear combination 
        $s = \sum_i w_i \phi_i$.
  \item The \emph{rate matrix} $\mathbf{K}$ of the coarse Markov model 
        built on discretised $s$ values.
\end{enumerate}
The spectral gap $\Delta\lambda = \lambda_1 - \lambda_2$ (where 
$\lambda_1=0$ and $\lambda_2$ is the second eigenvalue) quantifies 
timescale separation.  A large $\Delta\lambda$ means that the slow 
process is well-separated from fast fluctuations:
\begin{equation}
  \{w_i^*\} = \operatorname*{arg\,max}_{\{w_i\}} \Delta\lambda(s).
\end{equation}
The optimisation is performed using Maximum Caliber~\citep{Jaynes1980} 
to build the rate matrix from short unbiased (or reweighted biased) 
trajectories.  A multi-dimensional extension (MD-SGOOP)~\citep{Smith2018SGOOP} 
allows sequential construction of higher-dimensional RC sets.

SGOOP is particularly valuable as a \emph{diagnostic tool}: given any 
candidate CV (including those learned by neural networks), one can 
compute its associated spectral gap and rank it against alternatives, 
without running additional simulations.

\subsection{Markov State Model validation}

Markov State Models (MSMs)~\citep{Pande2010} provide a complementary 
kinetic quality assessment framework.  After projecting trajectories 
onto a candidate CV set and building a discrete-state MSM, quality is 
assessed by:
\begin{itemize}[noitemsep]
  \item \textbf{Implied timescales test}: Slow implied timescales should 
        plateau as the MSM lag time $\tau$ increases.  Non-convergence 
        indicates the CV is not Markovian at that resolution.
  \item \textbf{Chapman--Kolmogorov test}: The propagated MSM should 
        agree with direct dynamics at multiples of the lag time.
  \item \textbf{Number and quality of metastable states} (e.g.\ via 
        PCCA+): A good CV set should yield well-separated metastable 
        macrostates.
\end{itemize}

%%==================================================================
\section{Time-Correlation-Based Quality Metrics}
\label{sec:correlation}
%%==================================================================

\subsection{Autocorrelation function and decorrelation time}

The \emph{normalised autocorrelation function} (NACF) of a CV 
$s(t)$ is
\begin{equation}
  C_s(\tau) = \frac{\langle \delta s(t+\tau)\,\delta s(t)\rangle}
                   {\langle \delta s^2\rangle},
  \quad \delta s = s - \langle s \rangle.
\end{equation}
The decorrelation (or integrated autocorrelation) time is
\begin{equation}
  \tau_{\mathrm{int}} = \int_0^\infty C_s(\tau)\,\mathrm{d}\tau.
\end{equation}
A long $\tau_{\mathrm{int}}$ signals that the CV describes a slow 
mode---desirable for identifying the relevant process---but it does not 
by itself confirm that all slow modes are captured.  In the context of 
enhanced sampling, a good CV will have a rapidly decaying NACF 
\emph{after} the bias has been applied (indicating efficient exploration), 
while its unbiased NACF should decay slowly (confirming it genuinely 
encodes a rare event).

\subsection{Time-Lagged Independent Component Analysis (TICA)}
\label{sec:tica}

TICA~\citep{Molgedey1994,Perez2013TICA} finds the linear combinations of 
input features that maximise time-lagged autocorrelation:
\begin{equation}
  \max_{\mathbf{w}} \; \frac{\mathbf{w}^\top \mathbf{C}(\tau) \mathbf{w}}
                             {\mathbf{w}^\top \mathbf{C}(0) \mathbf{w}},
\end{equation}
where $\mathbf{C}(\tau)_{ij} = \langle \phi_i(t)\,\phi_j(t+\tau)\rangle$ 
is the time-lagged cross-correlation matrix and $\mathbf{C}(0)$ is the 
zero-lag covariance.  The leading TICA components (TICs) are the slowest 
linear modes of the system and serve directly as CVs.

\textbf{Quality criterion}: The TICA eigenvalue $\mu_k$ of the $k$-th 
TIC is related to its implied timescale by 
$t_k = -\tau/\ln|\mu_k|$.  A large $t_k$ indicates that TIC $k$ 
encodes a genuinely slow mode.

\textbf{Recent development---tICA-metadynamics}~\citep{Smith2024tICAMeta}: 
TICA-derived CVs are used directly in metadynamics, achieving $\sim$1.5$\times$ 
faster convergence of free energy surfaces and correct recovery of 
kinetics compared to intuition-based CVs.  Care must be taken to 
correct for periodic boundary artefacts in translational CVs and to 
apply trajectory reweighting when constructing TICA from biased data.

\subsection{Deep-TICA}

Deep-TICA~\citep{DeepTICA2021} extends TICA using neural networks as 
basis functions, enabling identification of \emph{non-linear} slow 
modes:
\begin{equation}
  \max_{\boldsymbol{\theta}} \; 
  \frac{\mathbf{f}_{\boldsymbol\theta}^\top \mathbf{C}(\tau) 
        \mathbf{f}_{\boldsymbol\theta}}
       {\mathbf{f}_{\boldsymbol\theta}^\top \mathbf{C}(0) 
        \mathbf{f}_{\boldsymbol\theta}},
\end{equation}
where $\mathbf{f}_{\boldsymbol\theta}(\mathbf{x})$ is the neural network 
output.  By maximising this objective (related to the Koopman operator 
approximation), Deep-TICA identifies the slowest non-linear dynamical 
modes.  This is particularly valuable when limited, non-equilibrium 
trajectories are available, as demonstrated in recent work on slow 
dynamics from short simulations~\citep{SlowModes2025}.

%%==================================================================
\section{Variational Approaches and the VAMP Score}
\label{sec:vamp}
%%==================================================================

\subsection{The variational approach to conformational dynamics (VAC/VAMP)}

The Variational Approach for Conformational dynamics (VAC) and its 
non-reversible generalisation VAMP~\citep{Wu2020VAMP} provide a 
rigorous upper-bound principle: \emph{any} approximation of the 
dominant Koopman eigenfunctions yields a score that is bounded by the 
true value.  The VAMP-$r$ score is defined as
\begin{equation}
  \mathcal{V}_r(\mathbf{f}) = 
  \left\|\mathbf{C}(0)^{-1/2}\,\mathbf{C}(\tau)\,\mathbf{C}(\tau)^\top
         \,\mathbf{C}(0)^{-1/2}\right\|_{S_r},
\end{equation}
where $\|\cdot\|_{S_r}$ is the Schatten-$r$ norm and $\mathbf{f}$ is 
any set of feature functions.  Maximising VAMP-2 is equivalent to 
finding the best rank-$k$ approximation to the transfer operator.

\textbf{Quality criterion}: A higher VAMP-2 score indicates better 
capture of the slow dynamics.  Importantly, the score can be computed 
from data alone and used for cross-validation and hyperparameter 
selection without any reference simulation.

\subsection{VAMPnets}

VAMPnets~\citep{Mardt2018VAMPnets} couple the VAMP score to a deep 
neural network architecture, performing end-to-end optimisation of:
\begin{enumerate}[noitemsep]
  \item Feature transformation (encoder).
  \item Soft state assignment (clustering).
  \item MSM estimation.
\end{enumerate}
The VAMP-2 score is the training objective, so the network is directly 
optimised for kinetic quality.  Extensions include:
\begin{itemize}[noitemsep]
  \item \textbf{GraphVAMPnets}~\citep{GraphVAMPnets2023}: Replaces the 
        standard encoder with a graph neural network (GNN), enabling 
        permutation- and rotation-invariant CV discovery, applied 
        successfully to self-assembly kinetics.
  \item \textbf{Independent Markov Decomposition}~\citep{Hempel2022IMD}: 
        Decomposes large macromolecules into independent Markovian 
        subsystems, using VAMP scores to validate each subsystem's CV 
        independently.
\end{itemize}

%%==================================================================
\section{Free-Energy Convergence as a Quality Proxy}
\label{sec:freeenergy}
%%==================================================================

The most direct practical test of CV quality in metadynamics is whether 
the reconstructed free energy surface (FES) converges.

\subsection{Block analysis and running averages}

Splitting the trajectory into blocks and computing the FES 
independently for each block gives a convergence profile.  A 
well-chosen CV leads to rapid convergence and small block-to-block 
variance.  The difference between FES estimates at successive times,
\begin{equation}
  \varepsilon(t) = \left\|F(s,t) - F(s,t-\Delta t)\right\|,
\end{equation}
is a standard convergence diagnostic in well-tempered 
metadynamics~\citep{Barducci2008}.

\subsection{Mean Force Integration (MFI)}

The MFI approach~\citep{Bjola2024MFI} estimates the FES by 
integrating the mean force along the CV from multiple independent 
(potentially asynchronous) biased simulations:
\begin{equation}
  F(s) = -\int_{s_0}^{s} \langle f(s') \rangle\,\mathrm{d}s',
\end{equation}
where $\langle f(s)\rangle$ is the reweighted mean force at value $s$.  
The variance of the mean force provides an intrinsic error estimate, 
making MFI a robust convergence metric that is agnostic to the specific 
biasing protocol.

\subsection{Effective CVs via metadynamics of paths}

\citet{Muller2024EffCVs} propose learning CVs from path-metadynamics 
data and assessing them by the rate of FES convergence.  CVs that lead 
to smooth, rapidly converged FES with clear barriers are preferred over 
those yielding rugged or slowly-converging surfaces.

%%==================================================================
\section{Information-Theoretic Metrics}
\label{sec:infotheory}
%%==================================================================

\subsection{Mutual information}

The \emph{mutual information} between a candidate CV $s$ and a set 
of reference slow coordinates $\{\phi_k\}$ quantifies how much 
information $s$ encodes about the important processes:
\begin{equation}
  I(s;\phi_k) = \sum_{s,\phi_k} P(s,\phi_k)\ln\frac{P(s,\phi_k)}{P(s)P(\phi_k)}.
\end{equation}
A good CV should have high mutual information with the true slow 
modes.  In practice---where the ``true'' slow coordinates are unknown---
mutual information can be estimated between pairs of candidate CVs 
or between CVs and experimental observables, and used for 
\emph{feature selection} prior to enhanced sampling.

\subsection{Transfer entropy}

Transfer entropy~\citep{Schreiber2000} is an asymmetric, 
information-theoretic measure of \emph{directed} influence:
\begin{equation}
  T_{Y\to X} = \sum P(x_{t+\tau},x_{t},y_{t})\ln
    \frac{P(x_{t+\tau}\mid x_{t},y_{t})}{P(x_{t+\tau}\mid x_{t})}.
\end{equation}
In the CV quality context, a large transfer entropy from a candidate 
CV to a known structural marker indicates that the CV genuinely 
\emph{drives} the structural change---it is causal, not merely 
correlated.

\subsection{Dimensionality and redundancy}

Information theory also provides metrics for redundancy among CVs.  
Given a set $\{s_1,\ldots,s_n\}$, the joint mutual information 
$I(\{s_i\};\phi_k)$ should increase sub-additively; CVs with high 
pairwise redundancy (high $I(s_i;s_j)$) can be pruned without loss.  
This principle underpins feature-selection methods used before TICA 
or VAMP analysis.

%%==================================================================
\section{Machine-Learning Frameworks Incorporating Quality Metrics}
\label{sec:ml}
%%==================================================================

\subsection{RAVE: Reweighted Autoencoded Variational Bayes}

RAVE~\citep{Tiwary2018RAVE} trains a variational autoencoder (VAE) 
on biased trajectory data, with an iterative protocol:
\begin{enumerate}[noitemsep]
  \item Run a short biased simulation.
  \item Train the VAE to reconstruct configurations and minimise the 
        KL divergence of the latent space.
  \item The low-dimensional latent variable becomes the new CV.
  \item Reweight trajectory data to remove the bias before 
        the next training step.
\end{enumerate}
Quality is assessed implicitly by the VAE reconstruction error and 
explicitly by running KS tests on extracted transition times.

\subsection{Deep Locally Non-linear Embedding (DeepLNE)}

DeepLNE~\citep{DeepLNE2024} constructs path-like CVs by combining a 
differentiable autoencoder with a $k$-nearest-neighbour graph:
\begin{equation}
  s(\mathbf{x}) = \frac{\sum_k w_k \lambda_k}{\sum_k w_k},
\end{equation}
where $\lambda_k$ are progress-along-path parameters derived from 
locally linear projections.  Quality is assessed by comparing the 
reconstructed free energy profile with reference calculations and by 
measuring how closely the CV tracks true transition pathways 
in benchmark systems (RNA tetraloop, alanine dipeptide).

\subsection{mlcolvar: A Unified Framework}

The \texttt{mlcolvar} library~\citep{Bonati2023mlcolvar} 
(available at \url{https://mlcolvar.readthedocs.io}) provides a 
modular, PLUMED-compatible Python framework implementing:
\begin{itemize}[noitemsep]
  \item Linear and non-linear dimensionality reduction 
        (PCA, autoencoders).
  \item Supervised CVs (DeepLDA, deep discriminant analysis).
  \item Dynamical CVs (Deep-TICA).
  \item Committor learning.
\end{itemize}
Quality assessment is built into the workflow via 
cross-validation, loss curves, and diagnostic notebooks, making it 
practical to compare CV architectures on equal footing.

\subsection{Uncertainty-based CVs}

\citet{Zaverkin2025UncertCV} propose using the \emph{prediction 
uncertainty} of a machine-learning potential as a CV for enhanced 
sampling.  The quality of this approach is assessed by whether new 
chemically relevant configurations discovered under the bias reduce 
model uncertainty, providing a self-consistent quality criterion 
that is particularly useful for building robust ML potential datasets.

%%==================================================================
\section{Comparative Summary and Practical Recommendations}
\label{sec:summary}
%%==================================================================

Table~\ref{tab:summary} summarises the principal metrics reviewed 
in this report, along with their key characteristics.

\begin{longtable}{@{}p{3.2cm}p{3.0cm}p{3.0cm}p{3.2cm}@{}}
\toprule
\textbf{Method / Metric} & \textbf{Category} & 
\textbf{Key quantity} & \textbf{Computational cost} \\
\midrule
\endfirsthead
\multicolumn{4}{c}{\small\itshape (continued from previous page)}\\
\toprule
\textbf{Method / Metric} & \textbf{Category} & 
\textbf{Key quantity} & \textbf{Computational cost} \\
\midrule
\endhead
\midrule
\multicolumn{4}{r}{\small\itshape (continues on next page)}\\
\endfoot
\bottomrule
\endlastfoot

Committor histogram test & 
Mechanistic (gold standard) & 
Distribution of $\hat{q}$ at TS & 
High (many short trajectories) \\[3pt]

Infrequent metadynamics KS test & 
Kinetic & 
$P(t_i)$ vs.\ Poisson & 
Moderate--high \\[3pt]

SGOOP spectral gap & 
Kinetic & 
$\Delta\lambda = \lambda_1-\lambda_2$ & 
Low (post-processing) \\[3pt]

MSM implied timescales & 
Kinetic & 
$t_k(\tau)$ plateau & 
Moderate \\[3pt]

Autocorrelation / 
$\tau_{\mathrm{int}}$ & 
Time-correlation & 
$C_s(\tau)$, $\tau_{\mathrm{int}}$ & 
Low (from MD traj.) \\[3pt]

TICA eigenvalue $\mu_k$ & 
Time-correlation & 
$t_k = -\tau/\ln|\mu_k|$ & 
Low--moderate \\[3pt]

VAMP-2 score & 
Variational & 
$\mathcal{V}_2(\mathbf{f})$ & 
Low (from data) \\[3pt]

FES convergence / block analysis & 
Thermodynamic & 
$\varepsilon(t)$ & 
Low (post-processing) \\[3pt]

Mean Force Integration & 
Thermodynamic & 
Variance of $\langle f(s)\rangle$ & 
Low--moderate \\[3pt]

Mutual information & 
Information-theoretic & 
$I(s;\phi_k)$ & 
Low--moderate \\[3pt]

Transfer entropy & 
Information-theoretic & 
$T_{Y\to X}$ & 
Moderate \\[3pt]

Committor-regularised ML & 
Mechanistic + ML & 
Committor loss term & 
Moderate (training) \\[3pt]

VAE / RAVE reconstruction & 
ML generative & 
Reconstruction error, KS test & 
Moderate (training + sim.) \\[3pt]

\end{longtable}
\captionof{table}{Comparative summary of CV quality metrics.}
\label{tab:summary}

\subsection{Are kinetic or time-correlation methods more common?}

Both families are widely used, but for different purposes:

\begin{itemize}[noitemsep]
  \item \textbf{Time-correlation methods} (TICA, VAMP score, Deep-TICA, 
        autocorrelation analysis) dominate as \emph{training objectives} 
        and \emph{offline selection criteria}, because they are 
        computationally cheap and can be evaluated from existing 
        trajectory data without additional simulations.

  \item \textbf{Kinetic methods} (infrequent metadynamics KS test, 
        SGOOP spectral gap, MSM implied timescales) are used as 
        \emph{validation checks} after an enhanced-sampling run, 
        providing a rigorous---if more costly---verification that 
        the extracted rates are physically meaningful.
\end{itemize}

In practice, the recommended workflow is:
\begin{enumerate}[noitemsep]
  \item Use TICA / VAMP score / Deep-TICA to select or train the CV.
  \item Validate with SGOOP spectral gap (low cost) and FES 
        convergence analysis.
  \item If kinetics are required, apply infrequent metadynamics with 
        the KS test.
  \item If a mechanistic interpretation is needed, perform committor 
        analysis on a subset of configurations at the transition state.
\end{enumerate}

\subsection{Special considerations for neural-network CVs}

When the CV is a neural network trained on sampling data (e.g.\ with 
RAVE, Deep-TICA, or mlcolvar), additional care is required:

\begin{enumerate}[noitemsep]
  \item \textbf{Cross-validation}: Split trajectory data into 
        training and test sets; monitor VAMP score or reconstruction 
        error on the test set to detect overfitting.
  \item \textbf{Reweighting}: When training on biased data, apply 
        appropriate reweighting (e.g.\ MBAR or metadynamics 
        reweighting) before computing TICA or VAMP scores.
  \item \textbf{Committor regularisation}: Penalise deviations from 
        committor behaviour during training to ensure the learned CV 
        reflects genuine transition state structure.
  \item \textbf{Explainability}: Use SHAP values or similar 
        tools to identify which input features the network focuses 
        on---this provides a sanity check on physical relevance.
\end{enumerate}

%%==================================================================
\section{Conclusions}
%%==================================================================

The assessment of CV quality in enhanced sampling simulations has 
matured considerably over the past decade.  The field now offers a 
rich toolkit spanning four complementary paradigms: 
(i)~kinetic assessment via spectral gap and barrier-crossing rates; 
(ii)~time-correlation analysis via TICA, Deep-TICA, and autocorrelation 
functions; (iii)~variational optimisation via VAMP scores and VAMPnets; 
and (iv)~mechanistic validation via the committor function.  
Information-theoretic measures (mutual information, transfer entropy) 
provide additional layer of validation, particularly for feature 
selection in high-dimensional systems.

Machine-learning frameworks such as RAVE, Deep-TICA, and mlcolvar 
incorporate these quality metrics directly into their training 
objectives, effectively automating CV optimisation and validation. 
The committor function remains the theoretical gold standard, and 
recent committor-regularised training strategies are beginning to 
make this standard practically achievable even for complex 
biomolecular systems.

For practitioners working with neural-network-learned CVs in 
metadynamics, the recommended minimum validation suite consists of: 
the VAMP-2 score (or Deep-TICA eigenvalue), the infrequent metadynamics 
KS test, and a spot-check committor histogram on configurations 
near the putative transition state.

%%==================================================================
%% References
%%==================================================================
\newpage
\begin{thebibliography}{99}

\bibitem[Laio and Parrinello(2002)]{Laio2002}
Laio, A.; Parrinello, M.
\newblock Escaping free-energy minima.
\newblock \emph{Proc.\ Natl.\ Acad.\ Sci.\ USA} \textbf{2002}, \emph{99}, 12562--12566.
\newblock \url{https://doi.org/10.1073/pnas.202427399}

\bibitem[Barducci et al.(2008)]{Barducci2008}
Barducci, A.; Bussi, G.; Parrinello, M.
\newblock Well-tempered metadynamics: A smoothly converging and tunable free-energy method.
\newblock \emph{Phys.\ Rev.\ Lett.} \textbf{2008}, \emph{100}, 020603.
\newblock \url{https://doi.org/10.1103/PhysRevLett.100.020603}

\bibitem[E and Vanden-Eijnden(2004)]{EWainanBerkovitz2004}
E, W.; Vanden-Eijnden, E.
\newblock Metastability, conformation dynamics, and transition pathways in complex systems.
\newblock In \emph{Multiscale Modelling and Simulation}, Springer, 2004; pp.~35--68.

\bibitem[Du et al.(1998)]{Du1998}
Du, R.; Pande, V.S.; Grosberg, A.Y.; Tanaka, T.; Shakhnovich, E.S.
\newblock On the transition coordinate for protein folding.
\newblock \emph{J.\ Chem.\ Phys.} \textbf{1998}, \emph{108}, 334--350.
\newblock \url{https://doi.org/10.1063/1.475393}

\bibitem[Geissler et al.(1999)]{Geissler1999}
Geissler, P.L.; Dellago, C.; Chandler, D.
\newblock Kinetic pathways of ion pair dissociation in water.
\newblock \emph{J.\ Phys.\ Chem.\ B} \textbf{1999}, \emph{103}, 3706--3710.
\newblock \url{https://doi.org/10.1021/jp984837g}

\bibitem[Tiwary and Berne(2016)]{Tiwary2016SGOOP}
Tiwary, P.; Berne, B.J.
\newblock Spectral gap optimization of order parameters for sampling complex molecular systems.
\newblock \emph{Proc.\ Natl.\ Acad.\ Sci.\ USA} \textbf{2016}, \emph{113}, 2839--2844.
\newblock \url{https://doi.org/10.1073/pnas.1600917113}

\bibitem[Smith et al.(2018)]{Smith2018SGOOP}
Smith, Z.; Pramanik, D.; Tsai, S.T.; Tiwary, P.
\newblock Multi-dimensional spectral gap optimization of order parameters (SGOOP) through conditional probability factorization.
\newblock \emph{J.\ Chem.\ Phys.} \textbf{2018}, \emph{149}, 234105.
\newblock \url{https://doi.org/10.1063/1.5064856}

\bibitem[Tiwary and Parrinello(2015)]{Tiwary2015IM}
Tiwary, P.; Parrinello, M.
\newblock From metadynamics to dynamics.
\newblock \emph{Phys.\ Rev.\ Lett.} \textbf{2013}, \emph{111}, 230602.
\newblock \url{https://doi.org/10.1103/PhysRevLett.111.230602}

\bibitem[Jaynes(1980)]{Jaynes1980}
Jaynes, E.T.
\newblock The minimum entropy production principle.
\newblock \emph{Annu.\ Rev.\ Phys.\ Chem.} \textbf{1980}, \emph{31}, 579--601.

\bibitem[Pande et al.(2010)]{Pande2010}
Pande, V.S.; Beauchamp, K.; Bowman, G.R.
\newblock Everything you wanted to know about Markov State Models but were afraid to ask.
\newblock \emph{Methods} \textbf{2010}, \emph{52}, 99--105.
\newblock \url{https://doi.org/10.1016/j.ymeth.2010.06.002}

\bibitem[Molgedey and Schuster(1994)]{Molgedey1994}
Molgedey, L.; Schuster, H.G.
\newblock Separation of a mixture of independent signals using time delayed correlations.
\newblock \emph{Phys.\ Rev.\ Lett.} \textbf{1994}, \emph{72}, 3634--3637.

\bibitem[P\'{e}rez-Hern\'{a}ndez et al.(2013)]{Perez2013TICA}
P\'{e}rez-Hern\'{a}ndez, G.; Paul, F.; Giorgino, T.; De Fabritiis, G.; No\'{e}, F.
\newblock Identification of slow molecular order parameters for Markov model construction.
\newblock \emph{J.\ Chem.\ Phys.} \textbf{2013}, \emph{139}, 015102.
\newblock \url{https://doi.org/10.1063/1.4811489}

\bibitem[Smith et al.(2024)]{Smith2024tICAMeta}
Smith, Z. et al.
\newblock tICA-metadynamics for identifying slow dynamics in membrane permeation.
\newblock \emph{J.\ Chem.\ Theory Comput.} \textbf{2024}.
\newblock \url{https://pmc.ncbi.nlm.nih.gov/articles/PMC11282584/}

\bibitem[Bonati et al.(2021)]{DeepTICA2021}
Bonati, L.; Piccini, G.; Parrinello, M.
\newblock Deep learning the slow modes for rare events sampling.
\newblock \emph{Proc.\ Natl.\ Acad.\ Sci.\ USA} \textbf{2021}, \emph{118}, e2113533118.
\newblock \url{https://doi.org/10.1073/pnas.2113533118}

\bibitem[Wu and No\'{e}(2020)]{Wu2020VAMP}
Wu, H.; No\'{e}, F.
\newblock Variational approach for learning Markov processes from time series data.
\newblock \emph{J.\ Nonlinear Sci.} \textbf{2020}, \emph{30}, 23--66.
\newblock \url{https://doi.org/10.1007/s00332-019-09567-y}

\bibitem[Mardt et al.(2018)]{Mardt2018VAMPnets}
Mardt, A.; Pasquali, L.; Wu, H.; No\'{e}, F.
\newblock VAMPnets for deep learning of molecular kinetics.
\newblock \emph{Nat.\ Commun.} \textbf{2018}, \emph{9}, 5.
\newblock \url{https://doi.org/10.1038/s41467-017-02388-1}

\bibitem[Huang et al.(2023)]{GraphVAMPnets2023}
Huang, X. et al.
\newblock GraphVAMPnets for uncovering slow collective variables of self-assembly dynamics.
\newblock \emph{J.\ Chem.\ Phys.} \textbf{2023}, \emph{159}, 094901.
\newblock \url{https://doi.org/10.1063/5.0158903}

\bibitem[Hempel et al.(2022)]{Hempel2022IMD}
Hempel, T. et al.
\newblock Deep learning to decompose macromolecules into independent Markovian domains.
\newblock \emph{Nat.\ Commun.} \textbf{2022}, \emph{13}, 7101.
\newblock \url{https://pmc.ncbi.nlm.nih.gov/articles/PMC9675806/}

\bibitem[Bjola et al.(2024)]{Bjola2024MFI}
Bjola, M. et al.
\newblock Estimating free-energy surfaces and their convergence from multiple biased simulations.
\newblock \emph{J.\ Chem.\ Theory Comput.} \textbf{2024}.
\newblock \url{https://doi.org/10.1021/acs.jctc.4c00091}

\bibitem[M\"{u}llender et al.(2024)]{Muller2024EffCVs}
M\"{u}llender, A. et al.
\newblock Effective data-driven collective variables for free energy calculations from metadynamics of paths.
\newblock \emph{PNAS Nexus} \textbf{2024}, \emph{3}, pgae159.
\newblock \url{https://doi.org/10.1093/pnasnexus/pgae159}

\bibitem[Schreiber(2000)]{Schreiber2000}
Schreiber, T.
\newblock Measuring information transfer.
\newblock \emph{Phys.\ Rev.\ Lett.} \textbf{2000}, \emph{85}, 461--464.
\newblock \url{https://doi.org/10.1103/PhysRevLett.85.461}

\bibitem[Tiwary and Berne(2018)]{Tiwary2018RAVE}
Tiwary, P.; Berne, B.J.
\newblock Predicting reaction coordinates in energy landscapes with diffuse metastable states.
\newblock \emph{Proc.\ Natl.\ Acad.\ Sci.\ USA} \textbf{2016}; 
\newblock See also: Sultan, M.M.; Wayment-Steele, H.K.; Pande, V.S.
\newblock Transferable neural networks for enhanced sampling of protein dynamics.
\newblock \emph{J.\ Chem.\ Theory Comput.} \textbf{2018}, \emph{14}, 1887--1894.

\bibitem[Chen and Tiwary(2024)]{DeepLNE2024}
Chen, M.; Tiwary, P.
\newblock Deep learning path-like collective variable for enhanced sampling.
\newblock \emph{J.\ Chem.\ Phys.} \textbf{2024}, \emph{160}, 174109.
\newblock \url{https://doi.org/10.1063/5.0202156}

\bibitem[Bonati et al.(2023)]{Bonati2023mlcolvar}
Bonati, L. et al.
\newblock A unified framework for machine learning collective variables for enhanced sampling.
\newblock \emph{J.\ Chem.\ Phys.} \textbf{2023}, \emph{159}, 014801.
\newblock \url{https://doi.org/10.1063/5.0156343}

\bibitem[Zaverkin et al.(2025)]{Zaverkin2025UncertCV}
Zaverkin, V. et al.
\newblock Enhanced sampling of robust molecular datasets with uncertainty-based collective variables.
\newblock \emph{J.\ Chem.\ Phys.} \textbf{2025}, \emph{162}, 034114.
\newblock \url{https://doi.org/10.1063/5.0238961}

\bibitem[Li et al.(2024)]{Li2024CCVSM}
Li, C. et al.
\newblock Committor-consistent variational string method.
\newblock \emph{J.\ Chem.\ Theory Comput.} \textbf{2024}.
\newblock \url{https://pmc.ncbi.nlm.nih.gov/articles/PMC11404425/}

\bibitem[Bonati et al.(2025)]{Bonati2025CommittorReg}
Bonati, L. et al.
\newblock Committor-regularized learning of differentiable collective variables.
\newblock \emph{J.\ Chem.\ Phys.} \textbf{2025}, \emph{164}, 104115.
\newblock \url{https://doi.org/10.1063/5.0252656}

\bibitem[Kang and Vanden-Eijnden(2024)]{Kang2024}
Kang, X.; Vanden-Eijnden, E.
\newblock Computing the committor with the committor to study the transition state.
\newblock \emph{Nat.\ Comput.\ Sci.} \textbf{2024}.
\newblock \url{https://zenodo.org/records/11164167}

\bibitem[Committor Flow(2025)]{CommittorFlow2025}
Authors, A. et al.
\newblock Following the committor flow: A data-driven discovery of transition mechanisms.
\newblock \emph{arXiv} \textbf{2025}, 2507.21961.
\newblock \url{https://arxiv.org/abs/2507.21961}

\bibitem[Slow Modes(2025)]{SlowModes2025}
Authors, B. et al.
\newblock Data efficient learning of molecular slow modes from nonequilibrium simulations.
\newblock \emph{ChemRxiv} \textbf{2025}.
\newblock \url{https://chemrxiv.org/doi/10.26434/chemrxiv-2025-vswkj}

\end{thebibliography}

\end{document}
